\section{Transactional Clean-Only Publication}
\label{sec:publication-semantics}

\subsection{Total outcome and result construction}

Acceptance becomes security-relevant only when it controls every externally
observable artifact. A pipeline that records a block but still returns candidate
bytes leaves the final decision to an error-prone caller. The public computation
is therefore total over tagged outcomes:
\begin{equation}
  \mathsf{Outcome}=\mathsf{Result}(r)\mid\mathsf{Error}(e).
  \label{eq:outcome-sum}
\end{equation}
A result $r=(v,b,h)$ contains a verdict $v\in\mathcal{V}$, an optional
deliverable byte vector $b$, and bounded diagnostics $h$. An error $e$ is typed
and carries no result object or output vector.

For a reconstruction candidate $c$, the result-construction relation is
\begin{equation}
\begin{aligned}
\mathsf{BuildResult}_{p}(a,c)&={}\\[-0.5ex]
&\begin{cases}
 (\Clean,\mathsf{Some}(c.y),h),
   & \exists k\ \mathsf{Acc}_{p}(a,c,k),\\
 (v,\mathsf{None},h),
   & \substack{v\in\{\Suspicious,\Blocked,\\
      \NotProcessed\}},\\
 \mathsf{Error}(e),
   & \text{a typed terminal failure occurs}.
\end{cases}
\end{aligned}
\label{eq:build-result}
\end{equation}
The first branch is the only branch that can contain deliverable bytes. In
particular,
\begin{equation}
\begin{aligned}
  \mathsf{Deliverable}(r)
  &\iff r.v=\Clean\\
  &\iff r.b=\mathsf{Some}(c.y),
  \qquad \text{for an accepted candidate }c.
\end{aligned}
  \label{eq:clean-only-delivery}
\end{equation}
The equality fixes the exposed bytes to the candidate that satisfied the full
acceptance predicate. A candidate size or digest can remain in bounded
diagnostics, but it does not identify a caller-accessible artifact.

Typed errors include unsafe input size or name, type disagreement, required
decoder failure, timeout, cancellation, panic containment, excessive output,
and persistence failure. Non-clean results and errors retain distinct diagnostic
meaning while sharing one publication consequence. Neither carries deliverable
bytes.

The invariant is repeated at interface boundaries. The native byte API clears
output for every non-clean verdict. The C ABI exposes zero output length and no
owned output allocation for non-clean results. The filesystem API proceeds only
from a clean result whose generated suffix and output profile have already
converged. Warning, bypass, dry-run, and unmodified fallback modes are outside
the accepted scientific relation.

\subsection{No-clobber filesystem transaction}

Let $o$ be the requested output path and $D=\mathsf{parent}(o)$. Input and output
that resolve to the same file are rejected. A pre-existing destination is never
removed as preparation for processing. A non-clean result or processing error
returns without touching it. A clean publication attempt reaches the
no-clobber operation, which returns a typed conflict and preserves the existing
bytes. This contract avoids converting a failed defense operation into data
loss.

For a clean result with bytes $y$, publication creates a uniquely named
temporary file $\tau$ in $D$, writes all of $y$, synchronizes the temporary
file, and performs a no-clobber persistence operation from $\tau$ to $o$.
Same-directory staging avoids a cross-filesystem move. If another actor creates
$o$ before persistence, the no-clobber operation fails and does not replace the
competing entry.

\begin{figure}[t]
\centering
\resizebox{0.98\textwidth}{!}{%
\begin{tikzpicture}[
  node distance=9mm and 12mm,
  state/.style={draw,rounded corners=2pt,align=center,minimum height=9mm,minimum width=24mm,font=\small},
  good/.style={state,fill=green!8},
  neutral/.style={state,fill=gray!12},
  fail/.style={state,fill=red!6},
  arrow/.style={-{Latex[length=2mm]},thick}
]
\node[neutral] (initial) {$S_0$: target absent};
\node[neutral,right=of initial] (staged) {$S_1$: temp complete};
\node[neutral,right=of staged] (synced) {$S_2$: temp synced};
\node[good,right=of synced] (published) {$S_3$: target visible};
\node[fail,below=of staged] (failed) {$S_F$: target absent or pre-existing target unchanged};
\draw[arrow] (initial) -- node[above,font=\scriptsize]{clean result only} (staged);
\draw[arrow] (staged) -- node[above,font=\scriptsize]{file sync} (synced);
\draw[arrow] (synced) -- node[above,font=\scriptsize]{persist no-clobber} (published);
\draw[arrow] (initial) -- node[left,font=\scriptsize]{non-clean or error} (failed);
\draw[arrow] (staged) -- (failed);
\draw[arrow] (synced) -- (failed);
\end{tikzpicture}%
}
\caption{Publication state machine. Only a clean validated candidate reaches
the same-directory staging path. Failure leaves the target absent or preserves
the pre-existing target observed outside this state sequence.}
\label{fig:publication-state-machine}
\end{figure}

For an outcome $u$, the publication transition is
\begin{equation}
  \mathsf{Publish}(o,u)=
  \begin{cases}
    \mathsf{commit}_{\mathrm{noclobber}}(o,y),
      & u=\mathsf{Result}(\Clean,\mathsf{Some}(y),h),\\
    \mathsf{no\mbox{-}op}(o),
      & u=\mathsf{Result}(v,\mathsf{None},h),\\
    \mathsf{no\mbox{-}op}(o),
      & u=\mathsf{Error}(e).
  \end{cases}
  \label{eq:publish-transition}
\end{equation}
The commit succeeds only after complete write, file synchronization, and the
no-clobber namespace operation. A write, sync, or persistence failure remains
an error. It is never converted into a clean published result.

\subsection{Publication invariants}

\begin{property}[No pre-commit visibility]
\label{prop:no-precommit-visibility}
Under the trusted-directory and atomic-namespace assumptions, an initially
absent deliverable path remains absent in $S_0$, $S_1$, $S_2$, and $S_F$. It
becomes visible only in $S_3$. A pre-existing target is unchanged on every
branch because no-clobber persistence cannot replace it.
\end{property}

The property concerns the designated output path. A process with directory
enumeration rights can observe a temporary name. Such observation is outside
the consumer publication interface and is constrained by directory permissions
and temporary-file creation semantics.

\begin{proposition}[Clean-only filesystem publication]
\label{prop:clean-only-filesystem}
If an invocation creates $o$ in state $S_3$, there is a candidate $c$ and
profile $k$ such that $c.y$ equals the target contents,
$R_{p,k}(x)\Downarrow c$, and $\mathsf{Acc}_{p}(a,c,k)$ holds.
\end{proposition}

\begin{proof}[Analytical argument]
The only incoming transition to $S_3$ is no-clobber persistence from $S_2$.
Staging is reachable only from the clean branch of
\cref{eq:build-result}, which requires an accepted candidate and exposes exactly
$c.y$. Complete write and file synchronization preserve that vector before the
namespace operation. The created target therefore contains the accepted
candidate, subject to the stated filesystem assumptions.
\end{proof}

\begin{proposition}[Failure preserves the destination]
\label{prop:failure-no-target}
A non-clean verdict, typed processing error, staging failure, synchronization
failure, or failed no-clobber persistence neither creates an initially absent
target nor modifies a pre-existing target.
\end{proposition}

\begin{proof}[Analytical argument]
Non-clean results and errors take no-op branches. Staging and synchronization
failures occur before the only namespace operation. The final operation has
no-clobber semantics, so conflict cannot replace an existing entry. Controlled
fault injection must still exercise these branches before the implementation
property is reported as empirically observed.
\end{proof}

A stale temporary file is not deliverable, although cleanup remains necessary
for confidentiality and storage hygiene. The transaction also distinguishes
atomic visibility from crash durability. The temporary file is synchronized
before persistence, but the current contract does not claim parent-directory
synchronization or survival across every power-loss and storage-controller
model.

No-clobber conflict reduces availability rather than integrity. The threat model
excludes an attacker with arbitrary write access to $D$, because such an actor
could create bytes at the deliverable path without invoking the defense process.

Publication does not prove that accepted media is harmless. It establishes that
the artifact exposed through the modeled interface is the same candidate that
satisfied source resolution, reconstruction level, complete profile parsing,
type convergence, signature policy, and resource bounds. Harmful semantic
content, a later renderer vulnerability, and unmodeled grammars remain outside
the claim.
